\bookchapter{第十三章：工具類App推薦}

\chapterintro{此章精選效率、健康、出行類工具，並補充實用系統與跨設備之能：清濁等清理與防騷擾、手機自帶攔截開關、天氣與地震預警；向日葵／ToDesk遠控電腦、Trae等AI辦公、人工桌面、LocalSend等局域網快傳、跨設備剪貼板、Aegis／Microsoft Authenticator動態口令、滴答清單日程與紙質掃描。}

\booksection{一、效率工具：管理汝之時與注意力}

效率工具乃數字生活之「骨架」，助爾將零散之想法、任務、資料整理得井井有條。下按場景分類推薦。

\booksubsection{待辦清單（任務管理）}

\begin{itemize}
  \item 滴答清單（TickTick）　所宜之人：須將工作、生活、學習任務與日程統一記錄者；此節重點推薦之一。其要：跨手機／電腦／平板／Web同步；清單+日曆+番茄鐘三合一；自然語言輸入「明天下午3點開會」「每週一交週報」自動識別時間與重複規則；支持子任務、標籤、優先級、附件與協作清單。宜「日程安排+任務打卡+每日回顧」一體化管理。
  \item 滴答清單如何用得更順： 收集箱：先悉數投入「今天／收集箱」，晚間再歸類至工作、家庭、學習清單。
  \item 日曆視圖：會議、截止日拖至日曆觀其衝突，較徒視待辦列表更直觀。
  \item 重複任務：服藥、繳費、健身設為規則任務，毋每次手填。
  \item 番茄鐘：於任務上徑啟專注，統計本週「果成幾何」。
  \item 與手機功能聯動：語音助手言「提醒我…」可入系統提醒；重要截止日兼啟系統日曆以為雙保險。
\end{itemize}

Todoist　所宜之人：求極簡、跨平臺之用戶。其要：界面克制清爽，項目層級清晰（任務子任務子子任務），karma積分激勵習慣養成。

Microsoft To Do　所宜之人：辦公族、Outlook用戶。其要：完全免費、無廣告，「我的一天」功能每日清空、減輕心理之負，與微軟生態深度集成。

系統便箋／提醒事項　所宜之人：不欲裝第三方App之輕度用戶。其要：原生體驗、零延遲，Siri／小愛語音徑直添加，與系統widget聯動。

\booksubsection{番茄鐘（專注力訓練）}

\begin{itemize}
  \item Forest　所宜之人：管不住手、易刷手機者。其要：專注時種樹，提前離App則樹枯，以「虛擬生命」建立正向反饋，可與好友合種。
  \item 潮汐　所宜之人：須白噪音輔專注／入睡者。其要：自然音效（雨聲、海浪、森林）+番茄鐘+冥想三合一，界面極簡療癒。
  \item 原則：經典25分鐘專注+5分鐘休息，可因自身狀態調為50/10或90/20。
\end{itemize}

\booksubsection{日曆（時間可視化）}

\begin{itemize}
  \item 系統日曆　所宜之人：所有用戶。其要：原生推送最穩、耗電極低，可與Siri／小愛語音添加。
  \item Google Calendar　所宜之人：多設備用戶、須與他人共享日程者。其要：「目標」功能自動為爾安排健身／學習時間，與Gmail智能識別郵件中之事件。
\end{itemize}

\booksubsection{筆記類（知識沉澱）}

\begin{itemize}
  \item Notion　所宜之人：須統一管理筆記、任務、數據庫、wiki者；學生、自由職業者、團隊協作。其要：全能知識庫，Block模塊化編輯，可搭建個人主頁、閱讀清單、項目看板，模板生態豐富。
  \item Obsidian　所宜之人：深度知識管理者、研究者、寫作者。其要：本地Markdown存儲，雙向鏈接+關係圖譜，使爾筆記成「第二大腦」，數據完全歸爾所有。
  \item flomo　所宜之人：欲記錄碎片靈感而不欲寫長文者。其要：卡片式筆記，如發朋友圈然記錄想法，支持\#標籤分類，按時間線回溯。
  \item 蘋果備忘錄　所宜之人：iPhone／Mac用戶、輕度記錄之需。其要：啟動快、支持掃描文檔／手繪／清單／iCloud同步，免費且與系統深度集成。
  \item 印象筆記　所宜之人：須大量網頁剪藏、資料歸檔之資深用戶。其要：網頁剪藏插件強大，OCR識別圖片文字，老牌穩定，然高級版須付費。
\end{itemize}

\booksubsection{PDF工具}

\begin{itemize}
  \item PDF Expert　所宜之人：iPad／iPhone用戶、須批註合同之商務人士、學生。其要：iOS上公認最強PDF編輯器，標註流暢、支持簽名、合併拆分、填寫表單。
  \item WPS　所宜之人：須同時處理Word／Excel／PPT／PDF之辦公族。其要：一站式辦公套件，國內網盤／微信文件兼容好，免費版功能已足日常。
  \item Adobe Acrobat　所宜之人：於PDF格式規範要求高者（法律、設計、出版）。其要：官方標準、轉換PDF/A格式無誤，跨設備同步與電子簽名功能專業。
\end{itemize}

\booksubsection{掃描類（紙質文件數字化）}

紙質合同、發票、作業、白板筆記、身份證複印件……手機攝像頭已能替代大部分掃描儀。核心流程：拍清楚自動裁邊與矯正導出PDF／圖片OCR隨選。

\begin{itemize}
  \item 掃描全能王（CamScanner）　所宜之人：屢須掃描合同、票據、證件之商務人士與學生。其要：邊緣自動識別、透視矯正、OCR文字識別，可生成多頁PDF並雲同步。
  \item 微軟Lens　所宜之人：學生、須將白板／書頁掃成Word者。其要：免費無廣告，可一鍵導出至OneNote／Word，識別文字準確率高。
  \item 蘋果備忘錄掃描　所宜之人：iPhone用戶、輕度掃描之需。其要：無需裝App，備忘錄內點相機圖標即可掃描，自動去陰影、簽名標註，免費好用。
  \item 白描／掃描全能王網頁版／品牌相機「文檔掃描」：華為、小米、OPPO、vivo系統相機大多內置「文檔／掃描」模式；安卓亦可用Google相機或系統掃一掃。宜應急拍一頁合同。
  \item 拍照小技：光線均勻、手機與紙面平行、對準四角；多頁文件按序拍畢再合併PDF；敏感證件留意水印與保存位置，用畢可自「最近刪除」徹底清理。
\end{itemize}

\booksection{二、健康運動}

健康類App乃手機最被低估之「私人助理」——能持續追蹤爾之睡眠、心率、運動量，較偶往醫院一次更能反映長期狀態。

\booksubsection{運動健身}

\begin{itemize}
  \item Keep　所宜之人：居家訓練者、零基礎新手、欲減脂塑形者。其要：海量免費課程，自HIIT至瑜伽全覆蓋，AI自動生成訓練計劃，社區氛圍濃厚。
  \item Nike Training Club（NTC）　所宜之人：求專業教練指導、望系統訓練者。其要：Nike官方出品，課程全部免費，動作演示由專業教練錄製，按肌群／器械分類清晰。
  \item 蘋果健身（Fitness）　所宜之人：Apple Watch用戶、欲以閉環激勵自己動起來者。其要：每日三活動圓環（移動／鍛煉／站立），目標可視化，與手表心率數據無縫打通。
\end{itemize}

\booksubsection{睡眠追蹤}

\begin{itemize}
  \item AutoSleep　所宜之人：Apple Watch用戶、欲知每晚深睡／淺睡比例者。其要：戴表入睡自動追蹤，無需手動啟停，予睡眠質量評分與就寢建議。
  \item Sleep Cycle　所宜之人：不戴表睡覺、欲以手機監測者。其要：以麥克風分析翻身／呼嚕聲，於淺睡階段智能喚醒，使爾起牀不苦。
  \item 小米運動健康　所宜之人：小米手環／手表用戶。其要：免費且與小米手環深度聯動，心率／血氧／睡眠一體追蹤，性價比高。
\end{itemize}

\booksubsection{冥想放鬆}

\begin{itemize}
  \item Calm　所宜之人：壓力大之上班族、入睡困難者。其要：「每日平靜」冥想引導+睡眠故事（名人朗讀），界面唯美，英文內容豐富。
  \item 潮汐　所宜之人：須白噪音專注／入睡之中文用戶。其要：國產精品，自然音效+冥想課程+番茄鐘，中文引導更親切。
  \item 小睡眠　所宜之人：須助眠音效、記錄夢話／呼嚕聲者。其要：海量白噪音可疊加組合（雨聲+風扇+篝火），睡前定時關閉，夢話錄音有趣又實用。
\end{itemize}

\booksubsection{女性健康}

\begin{itemize}
  \item Flo　所宜之人：須記錄經期、備孕或關注女性健康之用戶。其要：AI預測下次經期／排卵日準確度高，社區可匿名提問，界面設計精緻。
  \item 美柚　所宜之人：國內女性用戶、望看中文社區與科普者。其要：本地化內容豐富，按經期／備孕／孕期／育兒分階段提供對應知識。
\end{itemize}

\booksubsection{用藥提醒}

\begin{itemize}
  \item 用藥助手　所宜之人：須長期服藥之慢性病患者、須照顧老人用藥之子女。其要：掃碼識別藥品、自動建立用藥清單，按時彈窗提醒，記錄服藥情況可分享予家屬。
\end{itemize}

\booksubsection{數據中樞}

\begin{itemize}
  \item 蘋果健康／Google Fit　所宜之人：所有智能手機用戶。其要：匯聚來自手機、手環及第三方健康App之所有數據（步數、心率、睡眠、經期等），乃統一之「健康檔案中心」。
\end{itemize}

\booksection{三、旅行出行}

出門在外，數款趁手App能使爾少踩無數之坑。下按出行全流程推薦。

\booksubsection{航班高鐵}

\begin{itemize}
  \item 航旅縱橫　所宜之人：屢乘飛機之商務人士。其要：官方民航數據，航班動態／登機口變更／延誤預警最權威，可查機型與座位圖。
  \item 12306　所宜之人：所有國內火車出行者。其要：官方唯一可信購票渠道，支持候補購票（有人退票自動搶），學生票／團體票辦理方便。
  \item 飛常準　所宜之人：於航班準點率敏感、須接送機者。其要：航班延誤預測準確，可分享行程予家人實時查看爾之位置與到達時間。
\end{itemize}

\booksubsection{地圖導航}

\begin{itemize}
  \item 高德地圖　所宜之人：國內城市用戶、自駕／叫車族。其要：實時路況準確、避堵能力強，叫車／合乘一站搞定，公共交通換乘方案豐富。
  \item 百度地圖　所宜之人：於地點信息、室內導航有需之用戶。其要：全景地圖／AR實景導航，景區／商場室內地圖細緻，語音包豐富。
  \item Google Maps　所宜之人：出國旅行者。其要：全球覆蓋最全、本地評價真實，街景與離線地圖乃異國他鄉之救星。
\end{itemize}

\booksubsection{翻譯}

\begin{itemize}
  \item 谷歌翻譯　所宜之人：出國旅行者、跨語言學習者。其要：「對話」實時雙向翻譯、「相機」即時取詞翻譯乃旅行神器，支持離線語言包下載。
\end{itemize}

\booksubsection{匯率換算}

\begin{itemize}
  \item 極簡匯率　所宜之人：出國旅行、海淘、跨境購物者。其要：界面極致清爽，支持160+貨幣，匯率實時更新，可同時對比多種貨幣，長按數字速換算。
\end{itemize}

\booksubsection{天氣}

\begin{itemize}
  \item 系統天氣／中國天氣網　所宜之人：所有人，查預報與官方氣象預警之首選。其要：廠商天氣App與「中國天氣」等官方渠道對接氣象臺預警，暴雨、颱風、高溫等預警以橫幅／通知推送，權威性高。
  \item 彩雲天氣　所宜之人：須精準至分鐘級降雨預報者。其要：「未來兩小時是否下雨」準確度業內領先，降雨帶動態地圖直觀可見，雨前主動推送提醒。
  \item Carrot Weather　所宜之人：求趣味與豐富widget之天氣發燒友。其要：毒舌AI語音播報天氣、超多自定義widget桌面組件、雷達圖專業，付費版功能更全。
  \item 預警要點：天氣類App須允許通知方收得到預警；預警分藍、黃、橙、紅四級，紅色示災害性天氣風險最高，宜按提示調整出行。更多開啟方式見本章第七節「天氣預警／自然災害預警／地震預警」。
\end{itemize}

\booksubsection{旅行攻略}

\begin{itemize}
  \item 馬蜂窩　所宜之人：好作詳細攻略、自駕遊／深度遊者。其要：用戶原創遊記內容豐富，目的地玩法／路線／避坑指南詳盡，可下載離線攻略。
  \item 小紅書　所宜之人：欲看真實體驗、尋小眾打卡點／美食之年輕旅行者。其要：圖文短視頻形式直觀，「XX城市3天攻略」等內容更新快、時效性強，搜索體驗好。
\end{itemize}

\booksection{四、自動化工具啟蒙}

自動化乃「以一次配置換長期省事」之槓桿，能將爾每日重複數十次之操作壓縮為零。此節僅作概念啟蒙，具體玩法與詳細教程請見第二十三章：編程思維與自動化啟蒙。

\begin{itemize}
  \item iOS「快捷指令」：蘋果官方自動化工具，無需編程，可一鍵完成「到家開WiFi+亮度調低+播報天氣」等組合動作。
  \item 安卓自動化：系統級場景聯動（如小米「自動任務」、華為「情景智能」）上手最簡；Tasker功能強大然學習曲線陡峭，宜極客。
  \item 核心思路：先識別「每天都要重複做之操作」，再思「能否以一個觸發條件+多個動作代之」。詳細示例與進階玩法見第23章。
\end{itemize}

\booksection{五、瀏覽器與搜索引擎：爾之信息入口}

瀏覽器乃爾通往互聯網之窗，搜索引擎乃爾於信息海中導航之羅盤。善用之，能使爾少走許多彎路。

\booksubsection{瀏覽器工作之原理：}

爾於瀏覽器中輸入網址後，瀏覽器則：(1)向DNS服務器查詢域名對應之IP地址；(2)向該IP地址之服務器發送HTTP／HTTPS請求；(3)服務器返回網頁代碼（HTML、CSS、JavaScript）；(4)瀏覽器解析代碼並渲染為爾所見之頁面。

\booksubsection{主流瀏覽器推薦：}

\begin{itemize}
  \item 微軟Edge：基於Chromium內核，速度快、兼容好。內置Copilot AI助手，可總結網頁、翻譯、對話。支持跨設備同步，個人強烈推薦作為主力瀏覽器。
  \item Safari：蘋果設備默認瀏覽器，與系統深度集成，省電高效。隱私保護好，然擴展插件較少。
  \item Chrome：全球使用率最高，擴展豐富。然佔記憶體大，手機端相對耗電。
\end{itemize}

\booksubsection{搜索引擎之擇：}

\begin{itemize}
  \item Bing：微軟旗下搜索引擎，與Edge深度集成。AI摘要功能可徑直給出答案而非僅鏈接，搜索效率甚高。圖片搜索體驗亦佳。
  \item Yandex：俄羅斯最大之搜索引擎，圖片搜索與反向圖片搜索能力甚強。找相似圖片、識別物體時甚為好用。
  \item 百度：國內信息搜索之主力，然廣告較多，須一定辨別之能。
  \item Google：全球最精準之搜索引擎，然國內須特殊網絡方能使用。
\end{itemize}

\booksubsection{高效信息檢索指南：}

\begin{itemize}
  \item 善用AI搜索：Bing之AI摘要、Edge之Copilot可徑直回答問題，助爾從海量結果中提煉關鍵信息，較逐條點擊鏈接高效多矣。
  \item 識別虛假信息與廣告：搜索結果中標「廣告」或「贊助」之條目乃付費推廣，非自然搜索結果。留意甄別信息來源之權威性，優先擇官方網站、知名媒體之內容。
  \item 正確下載應用：於搜索引擎中輸入「應用名+官網」，如「微信官網」，優先點擊帶官方域名標識之結果。勿從第三方下載站安裝軟件，彼常捆綁流氓軟件。
  \item 官網鑑別之技：下載Windows 11應訪microsoft.com，下載Chrome應訪google.com。視域名與品牌一致與否，乃最基本之鑑別法。
\end{itemize}

\booksubsection{網站域名知識：}

域名者，網站之「門牌號」也，識之能助爾判斷網站性質與來源：

\begin{itemize}
  \item 常見後綴之義： .com — 商業機構，最常見
  \item .cn — 中國國家域名
  \item .org — 非營利組織
  \item .net — 網絡服務機構
  \item .edu — 教育機構
  \item .gov — 政府機構
  \item .xyz/.tech/.info — 新興通用域名，門檻低，須慎
\end{itemize}

國家／地區域名： .jp — 日本

.kr — 韓國

.uk — 英國

.de — 德國

.tw — 臺灣

.hk — 香港

安全提示：.gov與.edu後綴之網站通常更可信，以申請門檻高也。.xyz、.top等廉價域名易為釣魚網站所利用，訪時須多留心眼。

\booksubsection{廣告攔截：}

\begin{itemize}
  \item AdGuard　所宜之人：為各類彈窗、貼片廣告、應用內廣告所擾之所有用戶。其要：系統級廣告攔截（不限於瀏覽器），可過濾App內橫幅廣告與開屏追蹤，省流量、加速頁面加載，付費版支持全設備。
\end{itemize}

\booksubsection{隱私瀏覽器：}

\begin{itemize}
  \item Brave　所宜之人：重隱私、不欲為廣告追蹤之日常瀏覽者。其要：基於Chromium內核，開箱即用屏蔽廣告與追蹤器，自帶Tor隱身模式，速度較Chrome快、佔記憶體更低。
  \item DuckDuckGo　所宜之人：望「搜一個詞不被全網追着推廣告」之用戶。其要：自帶同名搜索引擎（不記錄搜索歷史），一鍵清除所有瀏覽痕跡（「燃燒按鈕」），跨平臺體驗統一。
\end{itemize}

\booksubsection{瀏覽器對比一覽：}

\begin{center}
\begin{tabularx}{\linewidth}{|X|X|X|}
\hline
\textbf{瀏覽器} & \textbf{其要} & \textbf{適用之人} \tabularnewline
\hline
微軟Edge & Chromium內核+內置Copilot AI，可總結網頁／翻譯／對話 & 求效率之辦公族、AI深度用戶（推薦主力） \tabularnewline
Safari & 與蘋果生態深度集成，省電高效，隱私保護強 & iPhone／Mac用戶，重續航與系統聯動 \tabularnewline
Chrome & 全球使用率最高，擴展插件生態最豐富 & 須大量擴展插件之進階用戶、跨設備同步黨 \tabularnewline
Brave & 開箱即用屏蔽廣告與追蹤器，速度更快 & 重隱私、不裝插件亦欲屏蔽廣告之用戶 \tabularnewline
DuckDuckGo & 同名搜索引擎不記錄歷史，一鍵清除痕跡 & 不欲被個性化追蹤之隱私敏感用戶 \tabularnewline
\hline
\end{tabularx}
\end{center}

\begin{tcbpractice}{試創建一個「睡前放鬆」之自動化流程}
\begin{enumerate}
  \item 創建觸發條件：特定時間（如22:30）。
  \item 添加執行動作：調低音量＞開啟勿擾模式＞播放放鬆歌單＞調低亮度。
  \item 保存並測試。
\end{enumerate}

\end{tcbpractice}

\booksection{六、實用系統工具：清理、防騷擾與安全防護}

效率與健康類App之外，實影響日常體驗者往往是「少被打擾、存儲夠用、騙子進不來」此三事。此節先盤點常用實用軟件，再單獨展開手機自帶之騷擾攔截功能——多人購機多年，卻從未開啟此等開關。

\booksubsection{系統清理與存儲優化}

\begin{itemize}
  \item 清濁　所宜之人：安卓用戶，欲輕量清理緩存、殘留檔案、空目錄，又不欲被「手機管家」全家桶打擾者。其要：體積小、無廣告之擾，支持深度掃描應用殘留、大檔案定位、存儲佔用分析；較商業「一鍵加速」更克制，宜動手能力強之用戶定期維護。留意：清理前先確認路徑，微信、相簿、文檔目錄勿盲目全選；系統級緩存盡量先用品牌自帶存儲清理，第三方工具僅作補充。
  \item 系統自帶存儲清理　所宜之人：所有人，尤長輩與怕折騰之用戶。其要：安卓［設定］＞［存儲］＞［清理加速］；iOS［設定］＞［通用］＞［iPhone存儲空間］。系統更諳自家App緩存結構，誤刪風險最低，宜為第一選擇。
  \item 微信／QQ內置存儲管理　所宜之人：聊天大戶。其要：微信［我］＞［設定］＞［通用］＞［存儲空間］，可按聊天、檔案類型精細清理；較外部「清微信」類工具更安全。
  \item 原則：能用系統／App自帶者，即不用第三方；能按路徑／類型勾選者，即勿「一鍵全清」。清理之目標在釋放空間，非將重要數據一併掃除。
\end{itemize}

\booksubsection{防騷擾與安全類實用軟件}

\begin{itemize}
  \item 國家反詐中心　所宜之人：所有人，尤長輩、財務相關崗位、屢網購／接陌生電話者。其要：官方出品，可開啟來電預警、短信風險提示、可疑網址檢測；與公安反詐數據庫聯動，乃「三道防線」中最有官方背書之一環。
  \item 品牌安全中心／手機管家　所宜之人：對應品牌用戶。其要：華為手機管家、小米安全中心、OPPO手機管家、vivo手機管家等，集成騷擾攔截、病毒查殺、權限管理、支付保護，通常免費且與系統底層協同更好，不建議再裝第二套商業安全軟件致權限相軋。
  \item 運營商防騷擾服務　所宜之人：被高頻外呼、AI語音電話轟炸之用戶。其要：移動發送KTFSR至10086開通；電信關注「天翼防騷擾」；聯通於「聯通手機管家」中設置。攔截發生於運營商側，手機不響鈴亦能擋住一批騷擾號。
  \item AdGuard　所宜之人：為各類彈窗、貼片廣告、應用內廣告所擾之所有用戶。其要：系統級廣告攔截（不限於瀏覽器），可過濾App內橫幅廣告與開屏追蹤，省流量、加速頁面加載，付費版支持全設備。
  \item 慎裝「清理加速／騷擾攔截」全家桶：部分第三方管家常駐後臺、反覆請求通訊錄與短信權限，攔截能力未必強於系統，反增耗電與隱私之險。先開系統自帶功能，再慮官方反詐App。
\end{itemize}

\textbf{實用軟件選型原則}\par

\begin{itemize}
  \item 系統優先：攔截、清理、權限管理，先開手機自帶開關。
  \item 官方補充：國家反詐中心、運營商防騷擾屬官方通道，可疊加。
  \item 第三方克制：同類功能裝一個即足，免權限相軋與後臺爭資源。
  \item 少即是多：工具裝得愈多，通知與權限之面愈廣，為「工具本身」所擾之概率愈高。
\end{itemize}

\booksubsection{手機自帶騷擾攔截功能：各品牌怎麼開}

下按品牌列出入口。名稱或隨系統版本微異，可於設置中徑搜「騷擾」「攔截」「未知來電」。與第八章「騷擾電話治理」互為補充：此處重開關位置，彼處重號碼標記與誤標申訴。

\begin{center}
\begin{tabularx}{\linewidth}{|X|X|X|}
\hline
\textbf{平臺} & \textbf{入口路徑} & \textbf{推薦開啟項} \tabularnewline
\hline
iPhone / iOS & ［設定］＞［電話］／［信息］ & 靜音未知來電；已阻止之聯絡人；信息中「過濾未知發件人」 \tabularnewline
小米／Redmi & 電話App＞設定＞騷擾攔截；或安全中心 & 智能攔截、陌生號響一聲掛斷、短信關鍵詞過濾 \tabularnewline
華為／榮耀 & 電話＞騷擾攔截；手機管家＞騷擾攔截 & 來電識別、黑名單、短信智能攔截 \tabularnewline
OPPO／realme／一加 & 電話App＞騷擾攔截；或本機＞騷擾攔截 & 來電攔截、詐騙電話識別、短信攔截規則 \tabularnewline
vivo / iQOO & 電話＞騷擾電話攔截；i管家 & 騷擾攔截、詐騙電話預警、信息過濾 \tabularnewline
三星 & 電話＞設定＞阻止號碼／來電顯示與垃圾電話 & 阻止未知號碼、垃圾電話識別（地區支持度不同） \tabularnewline
\hline
\end{tabularx}
\end{center}

\begin{tcbinfo}{三層防護疊加建議}
日常推薦「系統攔截+運營商防騷擾+國家反詐中心」三層疊加：系統管標記與黑名單，運營商於網側擋高頻外呼，反詐中心補官方預警。三層不必皆開至最嚴，免將快遞、外賣、醫生電話亦誤殺——可對通訊錄與最近聯絡人放行。

\end{tcbinfo}

\begin{tcbpractice}{動手練習：今日就把騷擾攔截開啟}
\begin{enumerate}
  \item 開電話App設定，尋得「騷擾攔截／靜音未知來電」，開啟智能攔截。
  \item 為爸媽手機亦檢查一遍：同路徑+開通運營商防騷擾（移動可發KTFSR至10086）。
  \item 應用商店安裝「國家反詐中心」，註冊並開啟來電預警。
  \item 試撥一個已被標記之測試號（或看攔截記錄），確認功能生效。
\end{enumerate}

\end{tcbpractice}

\booksection{七、天氣預警、自然災害預警與地震預警}

騷擾攔截所擋者乃「電話詐騙」，預警功能所擋者乃「天災與極端天氣」。多人只裝天氣App看溫度，卻未開通知、未開地震預警——暴雨紅色預警、颱風、地震來臨前數秒至數十秒，手機本可為爾爭得寶貴之避險時間。

\booksubsection{天氣預警：暴雨、颱風、高溫怎麼收}

\begin{itemize}
  \item 官方渠道優先：中國氣象局／中央氣象臺／各省氣象局發布之預警最權威。可關注「中國天氣」微信公眾號、微博，或下載「中國天氣」等官方類App。
  \item 系統天氣App：華為天氣、小米天氣、OPPO天氣、vivo天氣、蘋果「天氣」等多會推送當地氣象臺預警；入天氣App設定，確認「氣象預警／災害性天氣」通知已開啟。
  \item 彩雲天氣等分鐘級降雨：宜通勤、接送子女前判斷「這半小時要不要帶傘」；預警能力仍以官方氣象臺為準。
  \item 預警顏色之義（通用）： 藍色：注意防範，一般照常出行、關注更新。
  \item 黃色：加強防範，減少不必要之戶外活動。
  \item 橙色：嚴重，盡量留於室內，遠離危險區域。
  \item 紅色：特別嚴重，聽從停課、停工、避險指引，勿強闖積水、山洪、大風路段。
\end{itemize}

常見預警類型：暴雨、暴雪、颱風、高溫、寒潮、大風、雷電、冰雹、道路結冰、大霧、乾旱、地質災害風險。

\booksubsection{自然災害預警：非止「下雨」}

\begin{itemize}
  \item 發布主體：氣象類由氣象臺負責；地質災害、山洪等常由自然資源、水利、應急管理部門聯合發布；地方政府亦經應急廣播、短信、12379等渠道發布。
  \item 如何多渠道接收： 手機系統通知 + 天氣 App 預警。
  \item 微信關注當地「應急管理局」「氣象局」公眾號。
  \item 12379國家預警信息發布相關渠道（短信／服務，視地區開通情況）。
  \item 電視、應急廣播、村／社區大喇叭——山區、汛期勿只盯手機。
\end{itemize}

戶外與出行提醒：颱風天勿貼近廣告牌、勿行樹下；暴雨勿入地下車庫、隧道、下穿橋；山區旅遊遇「地質災害氣象風險」預警應立即離開陡坡、溝谷。

與急救手冊聯動：災後失聯、手機沒電、須定位求救時，見第三十二章地震／災害場景，以及第七章衛星通信與無網通。

\booksubsection{地震預警：與「地震預報」非一回事}

地震預警（EEW）者，於地震波（尤破壞力更強之橫波）到達目標區之前，利用電磁波快於地震波、以及臺站更快捕捉縱波之原理，向可能受影響區域提前數秒至數十秒推送警報。其不能預測地震是否發生，然能搶出「就近避險、切斷火源、衝出電梯」之窗口。距震中愈近，預警時間愈短，甚至可能來不及。

\begin{itemize}
  \item 官方與行業方案：中國大陸常見接入方包括中國地震預警網相關能力，以及成都高新減災研究所等機構推動之預警服務；部分廠商手機會預裝或可聯網接收預警。
  \item 第三方App：應用商店搜「地震預警」等官方／機構相關應用（安裝前看清開發者，避免山寨「假預警」App要爾開一堆無關權限）。
  \item 收到警報後如何做：室內就近躲避（桌下、承重牆角、堅固傢具旁），護住頭頸；遠離窗戶、外牆、高大傢具；勿盲目奔樓梯／跳樓；在電梯者即刻於最近樓層離開；室外到開闊地，遠離電線、廣告牌、危房。震後檢查燃氣與電路，留意餘震。
  \item 勿擾勿蔽：地震預警應允許「突破靜音／勿擾」響鈴；睡前置手機於伸手可及處，鈴聲音量勿調至0。
\end{itemize}

\booksubsection{手機自帶地震預警：各品牌怎麼開}

下列入口隨系統版本或微異，可於設置中徑搜「地震」「預警」「安全預警」。iOS系統本身對國內地震預警支持有限，通常須安裝可靠之第三方預警App，並保持定位與通知權限。

\begin{center}
\begin{tabularx}{\linewidth}{|X|X|X|}
\hline
\textbf{平臺} & \textbf{常見入口} & \textbf{說明} \tabularnewline
\hline
小米／Redmi & 手機管家 / 安全中心 > 地震預警；或 設置搜「地震」 & 部分版本系統內置預警服務，須聯網並開啟定位 \tabularnewline
華為／榮耀 & 設定＞安全＞應急預警通知／地震預警 & 名稱或為「應急預警通知」，內含地震等 \tabularnewline
OPPO／realme／一加 & 設定＞安全（SOS緊急聯絡）＞地震預警 & 部分機型視地區服務覆蓋而定 \tabularnewline
vivo / iQOO & 天氣App或設定＞安全／更多＞地震預警 & 以實際系統版本菜單為準 \tabularnewline
iPhone / iOS & 系統設置中通常無完整內置開關；安裝可靠地震預警App並允許通知 & App須後臺運行與定位權限；系統勿擾勿屏蔽其警報 \tabularnewline
三星等 & 設定＞安全與緊急情況／無線緊急警報 & 「緊急警報」覆蓋範圍因地區與運營商而異 \tabularnewline
\hline
\end{tabularx}
\end{center}

\begin{tcbwarn}{關於地震預警之常見誤區}
\begin{itemize}
  \item 預警$\neq$預報：其不能告爾「明天會不會地震」，只在地震已發、波尚未到達時報警。
  \item 非處處皆能提前多秒：震中附近可能僅數秒甚至來不及，預警主要惠及周邊區域。
  \item 關通知即同未裝：再好之預警，靜音、殺後臺、不給定位，皆可能收不到。
  \item 山寨App：警惕要通訊錄、短信權限之「地震預報」軟件，優先系統內置或正規機構方案。
\end{itemize}

\end{tcbwarn}

\begin{tcbinfo}{預警能力如何搭配}
日常建議：「系統天氣預警通知+官方氣象／應急渠道+地震預警（系統或正規App）」。與防騷擾同，不必裝一堆同類軟件；關鍵在通知權限開啟、定位允許、勿擾能突破。為長輩手機做完騷擾攔截後，順手將天氣預警與地震預警亦檢查一遍。

\end{tcbinfo}

\begin{tcbpractice}{動手練習：10分鐘檢查預警設置}
\begin{enumerate}
  \item 開系統天氣App，入設定，開啟「氣象預警／災害性天氣」通知。
  \item 於系統設置搜「地震」或「預警」，尋得並開啟地震預警／應急預警（安卓）；iPhone則安裝正規地震預警App並允許通知。
  \item 關注當地氣象局或應急局公眾號／服務號，允許其推送。
  \item 確認預警聲音不會被勿擾完全靜音；為爸媽手機同樣檢查一遍。
\end{enumerate}

\end{tcbpractice}

\booksection{八、跨設備協同：遠控電腦、快傳與帳號安全}

手機用得順，多半因其與電腦、平板「結隊」。此節補充數類高頻實用工具：手機遠控電腦、AI辦公與遠程協作、動態壁紙、局域網快傳、跨設備剪貼板、動態口令2FA，以及紙質掃描與日程任務之銜接。更系統之多設備聯動見第十七章。

\booksubsection{手機遠控電腦：向日葵／ToDesk}

\begin{itemize}
  \item 向日葵遠程控制　所宜之人：須於外出時操作家中／公司Windows或Mac之用戶；亦常用於助長輩遠程修手機（見第二十七章）。其要：國產生態成熟，手機App控制電腦桌面、傳檔案、開關機（部分須硬件）；免費版一般足個人輕度使用。
  \item ToDesk　所宜之人：求連接穩定、低延遲遠程辦公／運維之用戶。其要：設備碼+密碼連接，手機控電腦、電腦控手機（視版本）、檔案傳輸與多顯示器切換較順手，國內節點體驗通常較好。
  \item 基本用法：電腦端安裝客戶端並登錄／保持在線手機端裝App以設備識別碼+連接密碼（或同一帳號）建立會話手機觸控板／手勢操作桌面。
  \item 安全紅線（務必記住）： 陌生人自稱「客服／公檢法／運營商」令爾裝向日葵、ToDesk並遠程爾之手機或電腦，一概拒之——此典型詐騙手法也（詳見第二十七章）。
  \item 自用之時：開啟二次驗證、修改默認密碼、不用時退出帳號或關閉無人值守；公司電腦遵守單位安全規定。
  \item 公共場合遠程辦公，留意周圍是否有人偷窺屏幕與口令。
\end{itemize}

\booksubsection{AI辦公軟件：Trae／WorkBuddy／Dumate}

\begin{itemize}
  \item Trae　所宜之人：編程與文檔協作用戶，中文環境友好。其要：字節跳動推出之AI IDE／編程助手，SOLO模式可一句話生成項目；日常寫代碼、改文檔效率高。若電腦端運行，可於移動端配合遠控或官方配套能力查看／跟進任務（以產品最新形態為準）。詳見第四十一章AI編程工具全景。
  \item WorkBuddy　所宜之人：望AI助整理會議紀要、郵件、匯報、辦公文檔之職場用戶。其要：面向辦公場景之AI助手，宜「起草—潤色—總結」此類高頻文案；部分工作流可於手機端發起，重活付電腦端完成。
  \item Dumate　所宜之人：須多人協作、文檔與知識管理一體化之團隊／個人。其要：AI協作向辦公工具，重資料匯聚與協同編輯；手機端可查看與輕量編輯，複雜排版建議於電腦大屏處理。
  \item 「手機端遠程控制電腦」何解：AI辦公能力多運行於電腦／網頁端；於手機經向日葵／ToDesk連回電腦，或用產品自身之手機App／網頁，則通勤時可「遙控」電腦上之AI任務。原則：重計算、長文檔付電腦；看進度、改數句付手機。
  \item 隱私提醒：上傳合同、身份證、客戶數據前看清產品隱私條款與企業合規要求；公司資料優先用單位指定工具。
\end{itemize}

\booksubsection{人工桌面：米哈遊動態壁紙（手機+PC）}

\begin{itemize}
  \item 人工桌面　所宜之人：喜米哈遊角色、欲免費動態壁紙之用戶。其要：米哈遊推出之免費動態壁紙應用，支持手機端與PC端；角色Live2D／動態演出可作桌面，部分機型／系統可設為鎖屏或桌面壁紙。
  \item 省電與性能：動態壁紙較靜態圖更耗電、佔GPU；低電量或舊機型可只在充電時使用，或降低幀率（若App提供選項）。
  \item 安裝來源：優先官網／正規應用商店，避免第三方「破解版」捆綁木馬；PC版留意安裝路徑與開機自啟是否需要。
\end{itemize}

\booksubsection{局域網文件快傳：LocalSend／飛鴿傳書／隔空投送}

\begin{itemize}
  \item LocalSend　所宜之人：安卓／Windows／iOS／Mac／Linux混用，欲跨平臺、不經公有雲快傳者。其要：開源、局域網直連，同一Wi-Fi下互傳照片視頻文檔，無需數據線、不經廠商服務器，速度取決於路由器。
  \item 飛鴿傳書　所宜之人：辦公室Windows電腦之間、舊局域網辦公場景。其要：老牌局域網傳輸／聊天工具，宜內網傳大檔；留意從正規渠道下載，企業內網遵守IT策略。
  \item 隔空投送（AirDrop）　所宜之人：iPhone、iPad、Mac用戶。其要：蘋果生態一鍵分享，藍牙發現+Wi-Fi直連傳輸，免流量；可設「僅聯絡人」以降陌生人騷擾。
  \item 其他補充：小米互傳、華為分享、vivo互傳等品牌方案於同品牌設備間甚快；微信「文件傳輸助手」方便然受限於文件大小與登錄設備，大文件更推薦LocalSend／隔空投送。
  \item 使用提示：雙方連同一Wi-Fi；公共Wi-Fi傳敏感文件有風險，優先私人熱點；傳輸完成後於接收方確認文件是否完整。
\end{itemize}

\booksubsection{跨設備剪貼板同步：KDE Connect／手機電腦互傳}

\begin{itemize}
  \item 蘋果「通用剪貼板」　所宜之人：iPhone+Mac用戶。其要：登錄同一Apple ID、開接力後，手機複製Mac徑直粘貼（反之亦然），文字與部分文件皆可。
  \item KDE Connect　所宜之人：安卓+Windows／Linux／Mac，欲開源剪貼板同步、通知鏡像、遠程輸入之用戶。其要：局域網配對後同步剪貼板、推送手機通知至電腦、以手機為幻燈片翻頁器／觸控板；數據不經第三方帳號。
  \item 品牌「手機電腦互傳」：華為多屏協同／超級終端、小米妙享、vivo互傳、OPPO互連、榮耀互聯等，常自帶跨屏剪貼板與文件拖曳，同品牌體驗最佳。
  \item Windows+安卓官方通道：「手機連接／Link to Windows」（Your Phone演進）可於電腦看手機通知、照片，並於支持之場景複製粘貼；三星等機型集成更深。
  \item 安全提醒：剪貼板可能殘留密碼、驗證碼；共用電腦上慎用全局同步，用畢於密碼框避免二次粘貼洩露；部分輸入法亦讀剪貼板，留意權限。
\end{itemize}

\booksubsection{動態口令2FA：Aegis／Microsoft Authenticator}

\begin{itemize}
  \item 何以用認證器而非短信：短信驗證碼可能遭遇SIM補卡劫持；TOTP動態口令於本地按時間算法生成，無需依賴手機號信號，更抗釣魚與攔截（仍須防木馬與截屏）。
  \item Aegis Authenticator　所宜之人：安卓用戶，好開源、本地可控者。其要：開源免費，支持加密備份與導入導出，界面乾淨，可以指紋／人臉解鎖令牌列表。
  \item Microsoft Authenticator　所宜之人：用微軟帳號、Office 365／企業帳號之用戶；跨iOS／安卓。其要：推送批准登錄、密碼監控與帳號備份（視產品版本），與Windows／企業身份體系結合緊密。
  \item Google Authenticator／2FAS等：生態內亦常見；關鍵在有加密備份，換機不失令牌。
  \item 換機必為：舊機遷移或用服務中「備用碼／恢復密鑰」；備用碼打印於紙存兩處，勿只存同一部手機。
  \item 與密碼管理器：部分密碼庫亦可存TOTP，方便然風險更集中；高價值帳號建議「密碼管理器+獨立認證器」分工。詳見第十四章。
\end{itemize}

\booksubsection{推薦組合：一日之中怎麼串起來}

\textbf{實用工作流之例}\par

\begin{itemize}
  \item 晨：滴答清單看今日日程天氣／預警掃一眼電腦開Trae／文檔，手機置一旁。
  \item 晝：紙質材料用掃描類App轉PDF，經LocalSend／隔空投送發至電腦；長文於電腦用AI辦公工具起草。
  \item 外出：用ToDesk／向日葵連回電腦繼續任務；重要帳號登錄走Aegis／Microsoft Authenticator動態口令。
  \item 夜：滴答清單勾選完成項；須時人工桌面換壁紙放鬆一下。
\end{itemize}

\begin{tcbwarn}{遠控與帳號安全口訣}
\begin{itemize}
  \item 他人令爾裝遠控軟件掛斷電話，自尋官方客服。
  \item 驗證碼、動態口令勿告任何人，包括「快遞」「客服」。
  \item 公司電腦遠控與安裝AI軟件前，先問IT能不能用。
  \item 公共網絡傳文件、遠程登錄時格外警惕。
\end{itemize}

\end{tcbwarn}

\begin{tcbpractice}{動手練習：把「跨設備」鏈路跑通一次}
\begin{enumerate}
  \item 手機與電腦連同一Wi-Fi，用LocalSend（或隔空投送）互傳一張照片。
  \item 於滴答清單建一條「今晚8點：整理掃描件」任務，確認手機與電腦都能看到。
  \item 為郵箱或微信開啟2FA（若支持），用Aegis或Microsoft Authenticator綁定，並保存備用碼。
  \item 於自己電腦安裝向日葵或ToDesk，用自己手機成功連一次，再退出或加固密碼，熟悉流程但不留安全後門。
\end{enumerate}

\end{tcbpractice}
